\chapter{Gaussian second variation and symmetry modes}\label{sec:shrinker-variation}
The shrinker equation is the Euler--Lagrange equation of Gaussian area, but its Hessian differs from the ordinary minimal-surface Jacobi form. We use the scalar convention $H=\operatorname{div}_\Sigma\nu$ and $\Phi=H-\langle x,\nu\rangle/2$. Thus $\Phi=0$ is Definition \ref{def:self-shrinker}. Sources are \cite{CM12}, Sections 4--5, and \cite{CCMS20}, Section 3.1. The notation $[u]=\int u\rho\,dA$ and $\mathcal D$ is that of Definition \ref{def:gaussian-drift}.

\begin{proposition}[Compact variations of noncompact Gaussian area]\label{prop:gaussian-compact-variation}
Let $\Sigma$ be a smooth two-sided boundaryless immersed hypersurface with $F(\Sigma)=[1]<\infty$. For a smooth variation supported on a compact subset of its parameter manifold, with initial normal velocity $f\nu$,
\[
 F'(0)=[f\Phi].
\]
Consequently it is critical for all such variations exactly when it is a self-shrinker.
\end{proposition}
\begin{proof}
The changing part of the integral is compact. There $dA'=fH\,dA$ and $\rho'=-(f/2)\langle x,\nu\rangle\rho$, by Proposition \ref{prop:normal-variation-formulas}. The product rule proves the formula. Tests with arbitrary compactly supported $f$ imply $\Phi=0$ by continuity. No boundary term at infinity is taken; the integrand is unchanged there. This extends Proposition \ref{prop:gaussian-first-variation} to the stated noncompact setting.
\end{proof}

\begin{definition}[The Gaussian Jacobi operator]\label{def:shrinker-jacobi}
For a smooth two-sided self-shrinker, define
\[
 Lf=\mathcal Df+(|A|^2+\tfrac12)f,
 \qquad Q(f,g)=[\langle\nabla f,\nabla g\rangle-(|A|^2+\tfrac12)fg]
\]
for compactly supported smooth functions. Write $Q(f)=Q(f,f)$. The sign convention is that a negative value of $Q$ decreases centered Gaussian area to second order. See \cite{CM12}, (4.13), and \cite{CCMS20}, (3.1). The drift operator $\mathcal D$ and the Jacobi operator $L$ are different operators.
\end{definition}

\begin{proposition}[Linearizing the shrinker equation]\label{prop:shrinker-linearization}
For a normal variation $f\nu$ of a smooth two-sided hypersurface in $\mathbb R^{n+1}$,
\[
 \Phi'= -\left(\Delta-\tfrac12\langle x^T,\nabla\cdot\rangle+|A|^2+\tfrac12\right)f.
\]
At a self-shrinker this is $-Lf$. A tangential velocity $Y$ adds $Y(\Phi)$, which vanishes at a self-shrinker.
\end{proposition}
\begin{proof}
Proposition \ref{prop:mean-curvature-linearization} gives $H'=-\Delta f-|A|^2f$. Since $\nu'=-\nabla f$,
\[
 (\langle x,\nu\rangle)'=f-\langle x^T,\nabla f\rangle.
\]
Subtract half this identity from the first one. Tangential variations differentiate the scalar function by reparametrization. This also explains the signs in \cite{CM12}, (4.10).
\end{proof}

\begin{theorem}[Second variation with fixed center and scale]\label{thm:shrinker-fixed-hessian}
Let $\Sigma$ be a smooth two-sided boundaryless immersed self-shrinker with finite Gaussian area. For every smooth compactly supported variation with initial normal velocity $f\nu$,
\[
 F''(0)=-[fLf]=Q(f).
\]
Moreover $-[fLg]=Q(f,g)=-[gLf]$ for compactly supported smooth $f,g$. This is the fixed-center case of \cite{CM12}, Theorem 4.14, and \cite{CCMS20}, Section 3.1.
\end{theorem}
\begin{proof}
Differentiate Proposition \ref{prop:gaussian-compact-variation}. At a shrinker the factors multiplying $\Phi$ vanish, including any acceleration contribution, so the derivative is $[f\Phi']=-[fLf]$. Proposition \ref{prop:shrinker-linearization} gives the operator; Lemma \ref{lem:gaussian-integration-parts} gives $Q$ and its bilinear symmetry. This is formal self-adjointness on compact tests; it does not specify a self-adjoint domain or a discrete spectrum on an arbitrary noncompact surface.
\end{proof}

\begin{proposition}[Translation and dilation modes]\label{prop:shrinker-translation-dilation}
On a smooth two-sided self-shrinker, for each constant vector $v$ put $a_v=\langle v,\nu\rangle$. Then
\[
 La_v=\tfrac12a_v,\qquad LH=H.
\]
These pointwise identities do not require compactness. A mode is an eigenfunction only when it is nonzero and belongs to the domain under consideration. Compare \cite{CM12}, Theorem 5.2.
\end{proposition}
\begin{proof}
Translate the immersion to $x+sv$. Its scalar curvature and normal are unchanged, so its residual is $\Phi_s=-s a_v/2$. The normal velocity is $a_v$; its tangential component contributes zero because $\Phi=0$. Linearization therefore yields $-La_v=-a_v/2$.

For the dilation $x_s=(1+s)x$, scalar curvature is $H/(1+s)$ and the normal is unchanged. Thus $\Phi_s=H((1+s)^{-1}-(1+s))$, with derivative $-2H$. The normal velocity is $\langle x,\nu\rangle=2H$. Linearization gives $-L(2H)=-2H$. These arguments are local calculations and do not assert convergence of noncompact integrals.
\end{proof}

\begin{proposition}[Rotations give Jacobi fields]\label{prop:shrinker-rotation-mode}
If $K$ is a skew-symmetric linear map of $\mathbb R^{n+1}$ and $\Sigma$ is a smooth two-sided self-shrinker, then
\[
 L\langle Kx,\nu\rangle=0.
\]
\end{proposition}
\begin{proof}
The orthogonal maps $e^{sK}$ preserve both mean curvature and the radial Gaussian weight. Thus $e^{sK}\Sigma$ has identically zero shrinker residual for every $s$. Its normal velocity is $\langle Kx,\nu\rangle$. Proposition \ref{prop:shrinker-linearization} gives the conclusion. If the rotation preserves $\Sigma$, this function can be identically zero. This construction gives an integrable symmetry field; it does not say that every element of the kernel integrates to a family of shrinkers.
\end{proof}

\begin{corollary}[Orthogonality of the geometric modes on a closed shrinker]\label{cor:shrinker-mode-orthogonality}
On a smooth closed two-sided self-shrinker,
\[
 [H\langle v,\nu\rangle]=0
\]
for every constant vector $v$. Nonzero rotation modes are likewise orthogonal to $H$ and all translation modes.
\end{corollary}
\begin{proof}
All smooth functions are compactly supported on a closed manifold. The bilinear symmetry in Theorem \ref{thm:shrinker-fixed-hessian} gives $[uLw]=[wLu]$. Apply the pointwise equations of Propositions \ref{prop:shrinker-translation-dilation} and \ref{prop:shrinker-rotation-mode}. Distinct eigenvalues $1,1/2,0$ force the respective inner products to be zero.
\end{proof}

\begin{theorem}[Every finite-Gaussian-area proper shrinker is unstable at fixed center]\label{thm:shrinker-fixed-instability}
Let $\Sigma^n\subset\mathbb R^{n+1}$ be a nonempty smooth complete properly immersed two-sided self-shrinker without boundary and with $F(\Sigma)<\infty$. There is $f\in C_c^\infty(\Sigma)$ with $Q(f)<0$. Thus no such shrinker is stable for Gaussian area with its center and scale fixed. This explains the distinction preceding \cite{CM12}, Section 4.3.
\end{theorem}
\begin{proof}
Choose smooth radial cutoffs $0\leq\eta_R\leq1$ equal to one on $B_R$, zero outside $B_{2R}$, with $|\nabla_\Sigma\eta_R|\leq C/R$. They are smooth at zero by being constant there, and properness makes their support compact on the parameter manifold. Dropping the nonpositive curvature term gives
\[
 Q(\eta_R)\leq \frac{C^2}{R^2}[1]-\tfrac12[\eta_R^2].
\]
Dominated convergence yields $[\eta_R^2]\to[1]>0$. Hence the right side is negative for large $R$. This does not require a global integral bound for $|A|^2$: its contribution is finite on each compact support and only improves the inequality. The conclusion concerns fixed center and scale, unlike Definition \ref{def:shrinker-f-stability}.
\end{proof}
